\begin{table*}[t]
\centering
\small
\caption{CDC 2023 temporal-test cohort availability and fixed-model equal-budget results at the target-specific 10\% absolute referral budget.}
\label{tab:main_results}
\begin{tabular}{lrrrrrrrrr}
\toprule
Outcome & Full $N$ & Early $N$ (\%) & Events, $n$ (\%) & Event ceiling & Budget $B$ & Early capture & Full capture & $\Delta$ capture, pp (95\% CI) & Additional TP \\
\midrule
Preterm delivery & 3,480,382 & 2,601,060 (74.73\%) & 303,001 (8.71\%) & 71.83\% & 348,038 & 18.83\% & 21.41\% & +2.58 (2.47–2.70) & +7,821 \\
NICU admission & 3,476,584 & 2,597,274 (74.71\%) & 303,285 (8.72\%) & 70.62\% & 347,658 & 17.77\% & 20.50\% & +2.73 (2.62–2.84) & +8,271 \\
Low birth weight & 3,480,316 & 2,600,093 (74.71\%) & 243,776 (7.00\%) & 69.39\% & 348,032 & 20.56\% & 23.98\% & +3.43 (3.28–3.57) & +8,356 \\
\bottomrule
\end{tabular}
\vspace{2pt}
\begin{minipage}{0.99\textwidth}
\footnotesize
Early entry denotes prenatal-care initiation in months 1--3. Event ceiling is the fraction of all outcome events located in the early-entry cohort. Early and full policies selected the same target-specific absolute budget. Capture is population event capture. Confidence intervals are paired full-size record-level bootstrap percentile intervals from 500 replicates.
\end{minipage}
\end{table*}
